\documentclass[11pt,a4paper]{article}
\usepackage[utf8]{inputenc}
\usepackage[T1]{fontenc}
\usepackage[spanish,es-minimal,es-noshorthands]{babel}
\usepackage{lmodern}
\usepackage[sfdefault,lining]{FiraSans}
\usepackage[a4paper,top=2.3cm,bottom=2.2cm,left=1.8cm,right=1.8cm,headheight=1.3cm,headsep=0.4cm,footskip=1cm]{geometry}
\usepackage{amsmath,amssymb}
\usepackage{icomma}
\usepackage{graphicx}
\usepackage[table]{xcolor}
\usepackage{tikz}
\usepackage{fontawesome5}
\usepackage{ccicons}
\usepackage{fancyhdr}
\usepackage{lastpage}
\usepackage{tasks}
\usepackage{needspace}
\usepackage[most]{tcolorbox}
\usepackage{array}
\usepackage[hidelinks]{hyperref}
\AtBeginDocument{\hypersetup{pdftitle={Ficha de ejercicios de la unidad 7: Funciones. 3º ESO},pdfsubject={Matemáticas 3º ESO},pdfkeywords={funciones, imagen, antiimagen, tabla de valores, gráficas, dominio, recorrido, crecimiento, máximos y mínimos, continuidad}}}

\definecolor{navy}{HTML}{1B2748}
\definecolor{accent}{HTML}{E8603A}
\definecolor{teal}{HTML}{1F9E90}
\definecolor{soft}{HTML}{EEF1F8}
\definecolor{muted}{HTML}{5D6685}

\setlength{\parindent}{0pt}
\setlength{\parskip}{3pt}
\settasks{label=\alph*), label-format=\bfseries\color{navy}, label-width=1.3em, item-indent=1.8em, column-sep=0.6em, after-item-skip=5pt, before-skip=3pt, after-skip=4pt}

% ---------- capçalera i peu ----------
\newcommand{\llicencia}{\href{https://creativecommons.org/licenses/by-nc-sa/4.0/deed.es}{\ccbyncsa}\ \ Licencia Creative Commons Atribución-NoComercial-CompartirIgual 4.0 Internacional (CC BY-NC-SA 4.0)}
\pagestyle{fancy}
\fancyhf{}
\fancyhead[L]{\raisebox{-0.35\height}{\includegraphics[height=0.95cm]{logo-gva.png}}}
\fancyhead[C]{\small\color{muted}Ficha de ejercicios · Unidad 7: Funciones · 3\textsuperscript{o} ESO}
\fancyhead[R]{\raisebox{-0.35\height}{\includegraphics[height=0.95cm]{logo-vermellar.png}}}
\fancyfoot[L]{\scriptsize\color{muted}\llicencia}
\fancyfoot[R]{\small\color{muted}\thepage}
\renewcommand{\headrulewidth}{0.4pt}
\renewcommand{\headrule}{\color{muted!40}\hrule height \headrulewidth}
\fancypagestyle{primera}{\fancyhf{}\renewcommand{\headrulewidth}{0pt}
  \fancyfoot[L]{\scriptsize\color{muted}\llicencia}
  \fancyfoot[R]{\small\color{muted}\thepage}}

% ---------- nivells i exercicis ----------
\newcommand{\estrelles}[1]{{\color{accent}\ifcase#1\or\faStar\faStar[regular]\faStar[regular]\or\faStar\faStar\faStar[regular]\or\faStar\faStar\faStar\fi}}
\newcounter{exer}
\newcommand{\ex}[1]{\par\needspace{4\baselineskip}\medskip\refstepcounter{exer}\noindent{\color{navy}\bfseries\theexer.}\ \ {\small\estrelles{#1}}\ \ }

\newcommand{\seccio}[2]{\par\needspace{10\baselineskip}\bigskip
  \begin{tcolorbox}[colback=navy,colframe=navy,coltext=white,arc=2mm,boxrule=0pt,left=3mm,right=3mm,top=1.2mm,bottom=1.2mm]
  {\large\bfseries #1\quad #2}
  \end{tcolorbox}\vspace{-2pt}}

\newcolumntype{C}[1]{>{\centering\arraybackslash}p{#1}}
\usepackage{pifont}

% quadre de recordatori (només a la fitxa d'exercicis)
\newcommand{\recorda}[1]{\par\smallskip\begin{tcolorbox}[colback=accent!7,colframe=accent!55,arc=2mm,boxrule=0.5pt,left=2.5mm,right=2.5mm,top=1mm,bottom=1mm]\small{\color{accent}\faLightbulb[regular]}\ \textbf{Recuerda:}\ #1\end{tcolorbox}}
\newcommand{\res}[1]{#1}
\newcommand{\rest}[1]{#1}

\def\fmneed{14}
% =====================================================================
% Capçalera FentMates de les fitxes: un color per a cada curs.
% Els colors són els de la pàgina inicial (LEVELS a assets/js/materials.js);
% el logo ve de _eines/logo/pdf/ (python3 _eines/logo/genera.py).
%
% Ús, al final del preàmbul de la unitat:
%   % =====================================================================
% Capçalera FentMates de les fitxes: un color per a cada curs.
% Els colors són els de la pàgina inicial (LEVELS a assets/js/materials.js);
% el logo ve de _eines/logo/pdf/ (python3 _eines/logo/genera.py).
%
% Ús, al final del preàmbul de la unitat:
%   % =====================================================================
% Capçalera FentMates de les fitxes: un color per a cada curs.
% Els colors són els de la pàgina inicial (LEVELS a assets/js/materials.js);
% el logo ve de _eines/logo/pdf/ (python3 _eines/logo/genera.py).
%
% Ús, al final del preàmbul de la unitat:
%   \input{../comu/fentmates.tex}
%   \fmcurs{1eso}
% i, al començament del document:
%   \capcalera{Matemàtiques}{1r ESO}{Fitxa d'exercicis · Unitat 2}{Divisibilitat}
% Opcional abans de \input: \def\fmneed{14} (espai mínim abans de cada secció).
% =====================================================================
\makeatletter
\def\fm@logodir{../../logo/pdf}
\providecommand{\fmneed}{10}
\def\fm@cursnom{}

% clau del curs, color i color del text damunt del color
\newcommand{\fm@curs}[3]{\@namedef{fm@color@#1}{#2}\@namedef{fm@text@#1}{#3}}
\fm@curs{1eso}{E6B84A}{navy}
\fm@curs{2eso}{3B82F6}{white}
\fm@curs{3eso}{FF7A50}{white}
\fm@curs{4eso-a}{A78BFA}{white}
\fm@curs{4eso-b}{6D28D9}{white}
\fm@curs{1batx-mi}{9B1C1C}{white}
\fm@curs{1batx-ccss}{EF6B6B}{white}
\fm@curs{2batx-mii}{1E6B47}{white}
\fm@curs{2batx-ccss}{2BB3A3}{white}

\newcommand{\fmcurs}[1]{%
  \@ifundefined{fm@color@#1}{\PackageError{fentmates}{Curs desconegut: #1}{}}{}%
  \edef\fm@tmp{\noexpand\definecolor{curs}{HTML}{\@nameuse{fm@color@#1}}%
               \noexpand\colorlet{curstext}{\@nameuse{fm@text@#1}}}\fm@tmp
  \colorlet{cursfosc}{curs!72!black}%
  \def\fm@logo{\fm@logodir/vertical-#1.pdf}%
  \def\fm@marca{\fm@logodir/marca-#1.pdf}}

% ---------- capçalera de la primera pàgina ----------
% #1 matèria (petit), #2 curs, #3 tipus de document i unitat, #4 títol
\newcommand{\capcalera}[4]{%
  \gdef\fm@cursnom{#2}%
  \par\noindent
  \begin{tcolorbox}[colback=curs,colframe=curs,boxrule=0pt,arc=2.2mm,width=3.5cm,height=2.15cm,
      valign=center,left=3mm,right=2mm,top=1mm,bottom=1mm,nobeforeafter,box align=center]
    \color{curstext}{\footnotesize #1}\par\vspace{1pt}{\fontsize{19}{21}\selectfont\bfseries #2\par}%
  \end{tcolorbox}\hspace{5mm}%
  \begin{minipage}[c]{\dimexpr\linewidth-3.5cm-5mm-2.9cm\relax}
    {\color{cursfosc}\small\bfseries\MakeUppercase{#3}}\par\vspace{2pt}
    {\color{navy}\fontsize{20}{23}\selectfont\bfseries #4\par}
  \end{minipage}\hfill
  \begin{minipage}[c]{2.6cm}\raggedleft\includegraphics[width=2.6cm]{\fm@logo}\end{minipage}\par
  \vspace{2pt}}

% ---------- barres de secció del color del curs ----------
\renewcommand{\seccio}[2]{\par\needspace{\fmneed\baselineskip}\bigskip
  \begin{tcolorbox}[colback=curs,colframe=curs,coltext=curstext,arc=2mm,boxrule=0pt,left=3mm,right=3mm,top=1.2mm,bottom=1.2mm]
  {\large\bfseries #1\quad #2}
  \end{tcolorbox}\vspace{-2pt}}

% ---------- estrelles de dificultat del color del curs (a Batxillerat no se'n posen) ----------
\@ifundefined{estrelles}{}{%
  \renewcommand{\estrelles}[1]{{\color{curs}\ifcase#1\or\faStar\faStar[regular]\faStar[regular]\or\faStar\faStar\faStar[regular]\or\faStar\faStar\faStar\fi}}}

% ---------- capçalera i peu de les pàgines ----------
% primera pàgina: logos institucionals al peu
\fancypagestyle{primera}{\fancyhf{}\renewcommand{\headrulewidth}{0pt}
  \fancyfoot[L]{\raisebox{-0.32\height}{\includegraphics[height=0.78cm]{logo-gva.png}}\hspace{5mm}%
    \raisebox{-0.32\height}{\includegraphics[height=0.78cm]{logo-vermellar.png}}}
  \fancyfoot[R]{\scriptsize\color{muted}\href{https://creativecommons.org/licenses/by-nc-sa/4.0/deed.ca}{\ccbyncsa}\ \ CC BY-NC-SA 4.0 · fentmates.org\qquad\small\thepage}}
% la resta: el símbol a l'esquerra i el curs a la dreta, amb el títol al centre com abans
\fancyhead[L]{\raisebox{-0.3\height}{\includegraphics[height=0.62cm]{\fm@marca}}}
\fancyhead[R]{\small\bfseries\color{cursfosc}\fm@cursnom}
\renewcommand{\headrule}{\color{curs!45}\hrule height \headrulewidth}
\makeatother

%   \fmcurs{1eso}
% i, al començament del document:
%   \capcalera{Matemàtiques}{1r ESO}{Fitxa d'exercicis · Unitat 2}{Divisibilitat}
% Opcional abans de \input: \def\fmneed{14} (espai mínim abans de cada secció).
% =====================================================================
\makeatletter
\def\fm@logodir{../../logo/pdf}
\providecommand{\fmneed}{10}
\def\fm@cursnom{}

% clau del curs, color i color del text damunt del color
\newcommand{\fm@curs}[3]{\@namedef{fm@color@#1}{#2}\@namedef{fm@text@#1}{#3}}
\fm@curs{1eso}{E6B84A}{navy}
\fm@curs{2eso}{3B82F6}{white}
\fm@curs{3eso}{FF7A50}{white}
\fm@curs{4eso-a}{A78BFA}{white}
\fm@curs{4eso-b}{6D28D9}{white}
\fm@curs{1batx-mi}{9B1C1C}{white}
\fm@curs{1batx-ccss}{EF6B6B}{white}
\fm@curs{2batx-mii}{1E6B47}{white}
\fm@curs{2batx-ccss}{2BB3A3}{white}

\newcommand{\fmcurs}[1]{%
  \@ifundefined{fm@color@#1}{\PackageError{fentmates}{Curs desconegut: #1}{}}{}%
  \edef\fm@tmp{\noexpand\definecolor{curs}{HTML}{\@nameuse{fm@color@#1}}%
               \noexpand\colorlet{curstext}{\@nameuse{fm@text@#1}}}\fm@tmp
  \colorlet{cursfosc}{curs!72!black}%
  \def\fm@logo{\fm@logodir/vertical-#1.pdf}%
  \def\fm@marca{\fm@logodir/marca-#1.pdf}}

% ---------- capçalera de la primera pàgina ----------
% #1 matèria (petit), #2 curs, #3 tipus de document i unitat, #4 títol
\newcommand{\capcalera}[4]{%
  \gdef\fm@cursnom{#2}%
  \par\noindent
  \begin{tcolorbox}[colback=curs,colframe=curs,boxrule=0pt,arc=2.2mm,width=3.5cm,height=2.15cm,
      valign=center,left=3mm,right=2mm,top=1mm,bottom=1mm,nobeforeafter,box align=center]
    \color{curstext}{\footnotesize #1}\par\vspace{1pt}{\fontsize{19}{21}\selectfont\bfseries #2\par}%
  \end{tcolorbox}\hspace{5mm}%
  \begin{minipage}[c]{\dimexpr\linewidth-3.5cm-5mm-2.9cm\relax}
    {\color{cursfosc}\small\bfseries\MakeUppercase{#3}}\par\vspace{2pt}
    {\color{navy}\fontsize{20}{23}\selectfont\bfseries #4\par}
  \end{minipage}\hfill
  \begin{minipage}[c]{2.6cm}\raggedleft\includegraphics[width=2.6cm]{\fm@logo}\end{minipage}\par
  \vspace{2pt}}

% ---------- barres de secció del color del curs ----------
\renewcommand{\seccio}[2]{\par\needspace{\fmneed\baselineskip}\bigskip
  \begin{tcolorbox}[colback=curs,colframe=curs,coltext=curstext,arc=2mm,boxrule=0pt,left=3mm,right=3mm,top=1.2mm,bottom=1.2mm]
  {\large\bfseries #1\quad #2}
  \end{tcolorbox}\vspace{-2pt}}

% ---------- estrelles de dificultat del color del curs (a Batxillerat no se'n posen) ----------
\@ifundefined{estrelles}{}{%
  \renewcommand{\estrelles}[1]{{\color{curs}\ifcase#1\or\faStar\faStar[regular]\faStar[regular]\or\faStar\faStar\faStar[regular]\or\faStar\faStar\faStar\fi}}}

% ---------- capçalera i peu de les pàgines ----------
% primera pàgina: logos institucionals al peu
\fancypagestyle{primera}{\fancyhf{}\renewcommand{\headrulewidth}{0pt}
  \fancyfoot[L]{\raisebox{-0.32\height}{\includegraphics[height=0.78cm]{logo-gva.png}}\hspace{5mm}%
    \raisebox{-0.32\height}{\includegraphics[height=0.78cm]{logo-vermellar.png}}}
  \fancyfoot[R]{\scriptsize\color{muted}\href{https://creativecommons.org/licenses/by-nc-sa/4.0/deed.ca}{\ccbyncsa}\ \ CC BY-NC-SA 4.0 · fentmates.org\qquad\small\thepage}}
% la resta: el símbol a l'esquerra i el curs a la dreta, amb el títol al centre com abans
\fancyhead[L]{\raisebox{-0.3\height}{\includegraphics[height=0.62cm]{\fm@marca}}}
\fancyhead[R]{\small\bfseries\color{cursfosc}\fm@cursnom}
\renewcommand{\headrule}{\color{curs!45}\hrule height \headrulewidth}
\makeatother

%   \fmcurs{1eso}
% i, al començament del document:
%   \capcalera{Matemàtiques}{1r ESO}{Fitxa d'exercicis · Unitat 2}{Divisibilitat}
% Opcional abans de \input: \def\fmneed{14} (espai mínim abans de cada secció).
% =====================================================================
\makeatletter
\def\fm@logodir{../../logo/pdf}
\providecommand{\fmneed}{10}
\def\fm@cursnom{}

% clau del curs, color i color del text damunt del color
\newcommand{\fm@curs}[3]{\@namedef{fm@color@#1}{#2}\@namedef{fm@text@#1}{#3}}
\fm@curs{1eso}{E6B84A}{navy}
\fm@curs{2eso}{3B82F6}{white}
\fm@curs{3eso}{FF7A50}{white}
\fm@curs{4eso-a}{A78BFA}{white}
\fm@curs{4eso-b}{6D28D9}{white}
\fm@curs{1batx-mi}{9B1C1C}{white}
\fm@curs{1batx-ccss}{EF6B6B}{white}
\fm@curs{2batx-mii}{1E6B47}{white}
\fm@curs{2batx-ccss}{2BB3A3}{white}

\newcommand{\fmcurs}[1]{%
  \@ifundefined{fm@color@#1}{\PackageError{fentmates}{Curs desconegut: #1}{}}{}%
  \edef\fm@tmp{\noexpand\definecolor{curs}{HTML}{\@nameuse{fm@color@#1}}%
               \noexpand\colorlet{curstext}{\@nameuse{fm@text@#1}}}\fm@tmp
  \colorlet{cursfosc}{curs!72!black}%
  \def\fm@logo{\fm@logodir/vertical-#1.pdf}%
  \def\fm@marca{\fm@logodir/marca-#1.pdf}}

% ---------- capçalera de la primera pàgina ----------
% #1 matèria (petit), #2 curs, #3 tipus de document i unitat, #4 títol
\newcommand{\capcalera}[4]{%
  \gdef\fm@cursnom{#2}%
  \par\noindent
  \begin{tcolorbox}[colback=curs,colframe=curs,boxrule=0pt,arc=2.2mm,width=3.5cm,height=2.15cm,
      valign=center,left=3mm,right=2mm,top=1mm,bottom=1mm,nobeforeafter,box align=center]
    \color{curstext}{\footnotesize #1}\par\vspace{1pt}{\fontsize{19}{21}\selectfont\bfseries #2\par}%
  \end{tcolorbox}\hspace{5mm}%
  \begin{minipage}[c]{\dimexpr\linewidth-3.5cm-5mm-2.9cm\relax}
    {\color{cursfosc}\small\bfseries\MakeUppercase{#3}}\par\vspace{2pt}
    {\color{navy}\fontsize{20}{23}\selectfont\bfseries #4\par}
  \end{minipage}\hfill
  \begin{minipage}[c]{2.6cm}\raggedleft\includegraphics[width=2.6cm]{\fm@logo}\end{minipage}\par
  \vspace{2pt}}

% ---------- barres de secció del color del curs ----------
\renewcommand{\seccio}[2]{\par\needspace{\fmneed\baselineskip}\bigskip
  \begin{tcolorbox}[colback=curs,colframe=curs,coltext=curstext,arc=2mm,boxrule=0pt,left=3mm,right=3mm,top=1.2mm,bottom=1.2mm]
  {\large\bfseries #1\quad #2}
  \end{tcolorbox}\vspace{-2pt}}

% ---------- estrelles de dificultat del color del curs (a Batxillerat no se'n posen) ----------
\@ifundefined{estrelles}{}{%
  \renewcommand{\estrelles}[1]{{\color{curs}\ifcase#1\or\faStar\faStar[regular]\faStar[regular]\or\faStar\faStar\faStar[regular]\or\faStar\faStar\faStar\fi}}}

% ---------- capçalera i peu de les pàgines ----------
% primera pàgina: logos institucionals al peu
\fancypagestyle{primera}{\fancyhf{}\renewcommand{\headrulewidth}{0pt}
  \fancyfoot[L]{\raisebox{-0.32\height}{\includegraphics[height=0.78cm]{logo-gva.png}}\hspace{5mm}%
    \raisebox{-0.32\height}{\includegraphics[height=0.78cm]{logo-vermellar.png}}}
  \fancyfoot[R]{\scriptsize\color{muted}\href{https://creativecommons.org/licenses/by-nc-sa/4.0/deed.ca}{\ccbyncsa}\ \ CC BY-NC-SA 4.0 · fentmates.org\qquad\small\thepage}}
% la resta: el símbol a l'esquerra i el curs a la dreta, amb el títol al centre com abans
\fancyhead[L]{\raisebox{-0.3\height}{\includegraphics[height=0.62cm]{\fm@marca}}}
\fancyhead[R]{\small\bfseries\color{cursfosc}\fm@cursnom}
\renewcommand{\headrule}{\color{curs!45}\hrule height \headrulewidth}
\makeatother

\fmcurs{3eso}

\begin{document}
\thispagestyle{primera}

\capcalera{Matemáticas}{3\textsuperscript{o} ESO}{Ficha de ejercicios · Unidad 7}{Funciones}
\vspace{10pt}
Nombre y apellidos: \rule{6.6cm}{0.4pt}\hfill Grupo: \rule{1.6cm}{0.4pt}\hfill Fecha: \rule{2.2cm}{0.4pt}
\vspace{8pt}

\begin{tcolorbox}[colback=soft,colframe=soft,arc=2mm,left=3mm,right=3mm,top=1.5mm,bottom=1.5mm]
\small Ejercicios graduados por niveles de dificultad:\quad \estrelles{1}\ básico \qquad \estrelles{2}\ medio \qquad \estrelles{3}\ avanzado.\\[2pt] Resuélvelos en el cuaderno mostrando todos los pasos. Para dibujar gráficas, utiliza papel cuadriculado y elige una escala adecuada para cada eje.
\end{tcolorbox}

% =====================================================
\seccio{1}{Concepto de función}

\ex{1} Indica si cada relación es una función:\begin{tasks}(2)
  \task A cada número natural, su cuadrado.
  \task A cada persona, sus hermanos.
  \task A cada cantidad de litros de gasolina, el precio que pagamos.
  \task A cada $x>0$, sus raíces cuadradas.
  \task A cada alumno, la nota del último examen.
  \task A cada nota, los alumnos que la han sacado.
\end{tasks}\recorda{Una relación es una función si a \textbf{cada} valor de la variable independiente ($x$) le corresponde \textbf{un único} valor de la variable dependiente ($y$).}

\ex{1} En cada relación, ¿cuál es la variable independiente y cuál la dependiente?\begin{tasks}(1)
  \task La duración de una llamada y lo que cuesta.
  \task El número de obreros y el tiempo que tardan en acabar una obra.
  \task El lado de un triángulo equilátero y su perímetro.
\end{tasks}

\ex{1} ¿Cuáles de estas gráficas corresponden a una función?\par\smallskip\begin{tasks}(4)
  \task \begin{tikzpicture}[x=0.27cm,y=0.27cm,baseline=(current bounding box.center)]
\draw[muted!22,very thin,xstep=1,ystep=1] (-5,-4) grid (5,4);
\draw[-latex,muted] (-5,0)--(5.35,0) node[right,font=\scriptsize]{$x$};
\draw[-latex,muted] (0,-4)--(0,4.35) node[above,font=\scriptsize]{$y$};
\node[font=\tiny,muted,below] at (-4,0) {-4};
\node[font=\tiny,muted,below] at (-2,0) {-2};
\node[font=\tiny,muted,below] at (2,0) {2};
\node[font=\tiny,muted,below] at (4,0) {4};
\node[font=\tiny,muted,left] at (0,-4) {-4};
\node[font=\tiny,muted,left] at (0,-2) {-2};
\node[font=\tiny,muted,left] at (0,2) {2};
\begin{scope}\clip (-5,-4) rectangle (5,4);
\draw[accent,line width=1.1pt] (-4.0,2.0)--(-2.0,-2.0)--(1.0,3.0)--(4.0,1.0);
\end{scope}
\fill[accent] (-4.0,2.0) circle (2.2pt);
\fill[accent] (4.0,1.0) circle (2.2pt);
\end{tikzpicture}
  \task \begin{tikzpicture}[x=0.27cm,y=0.27cm,baseline=(current bounding box.center)]
\draw[muted!22,very thin,xstep=1,ystep=1] (-5,-4) grid (5,4);
\draw[-latex,muted] (-5,0)--(5.35,0) node[right,font=\scriptsize]{$x$};
\draw[-latex,muted] (0,-4)--(0,4.35) node[above,font=\scriptsize]{$y$};
\node[font=\tiny,muted,below] at (-4,0) {-4};
\node[font=\tiny,muted,below] at (-2,0) {-2};
\node[font=\tiny,muted,below] at (2,0) {2};
\node[font=\tiny,muted,below] at (4,0) {4};
\node[font=\tiny,muted,left] at (0,-4) {-4};
\node[font=\tiny,muted,left] at (0,-2) {-2};
\node[font=\tiny,muted,left] at (0,2) {2};
\begin{scope}\clip (-5,-4) rectangle (5,4);
\draw[teal,line width=1.1pt] plot[smooth] coordinates {(3.000,0.000) (2.984,0.314) (2.934,0.624) (2.853,0.927) (2.741,1.220) (2.598,1.500) (2.427,1.763) (2.229,2.007) (2.007,2.229) (1.763,2.427) (1.500,2.598) (1.220,2.741) (0.927,2.853) (0.624,2.934) (0.314,2.984) (0.000,3.000) (-0.314,2.984) (-0.624,2.934) (-0.927,2.853) (-1.220,2.741) (-1.500,2.598) (-1.763,2.427) (-2.007,2.229) (-2.229,2.007) (-2.427,1.763) (-2.598,1.500) (-2.741,1.220) (-2.853,0.927) (-2.934,0.624) (-2.984,0.314) (-3.000,0.000) (-2.984,-0.314) (-2.934,-0.624) (-2.853,-0.927) (-2.741,-1.220) (-2.598,-1.500) (-2.427,-1.763) (-2.229,-2.007) (-2.007,-2.229) (-1.763,-2.427) (-1.500,-2.598) (-1.220,-2.741) (-0.927,-2.853) (-0.624,-2.934) (-0.314,-2.984) (-0.000,-3.000) (0.314,-2.984) (0.624,-2.934) (0.927,-2.853) (1.220,-2.741) (1.500,-2.598) (1.763,-2.427) (2.007,-2.229) (2.229,-2.007) (2.427,-1.763) (2.598,-1.500) (2.741,-1.220) (2.853,-0.927) (2.934,-0.624) (2.984,-0.314) (3.000,-0.000)};
\end{scope}
\end{tikzpicture}
  \task \begin{tikzpicture}[x=0.27cm,y=0.27cm,baseline=(current bounding box.center)]
\draw[muted!22,very thin,xstep=1,ystep=1] (-5,-4) grid (5,4);
\draw[-latex,muted] (-5,0)--(5.35,0) node[right,font=\scriptsize]{$x$};
\draw[-latex,muted] (0,-4)--(0,4.35) node[above,font=\scriptsize]{$y$};
\node[font=\tiny,muted,below] at (-4,0) {-4};
\node[font=\tiny,muted,below] at (-2,0) {-2};
\node[font=\tiny,muted,below] at (2,0) {2};
\node[font=\tiny,muted,below] at (4,0) {4};
\node[font=\tiny,muted,left] at (0,-4) {-4};
\node[font=\tiny,muted,left] at (0,-2) {-2};
\node[font=\tiny,muted,left] at (0,2) {2};
\begin{scope}\clip (-5,-4) rectangle (5,4);
\draw[teal,line width=1.1pt] plot[smooth] coordinates {(5.000,-4.000) (4.605,-3.900) (4.220,-3.800) (3.845,-3.700) (3.480,-3.600) (3.125,-3.500) (2.780,-3.400) (2.445,-3.300) (2.120,-3.200) (1.805,-3.100) (1.500,-3.000) (1.205,-2.900) (0.920,-2.800) (0.645,-2.700) (0.380,-2.600) (0.125,-2.500) (-0.120,-2.400) (-0.355,-2.300) (-0.580,-2.200) (-0.795,-2.100) (-1.000,-2.000) (-1.195,-1.900) (-1.380,-1.800) (-1.555,-1.700) (-1.720,-1.600) (-1.875,-1.500) (-2.020,-1.400) (-2.155,-1.300) (-2.280,-1.200) (-2.395,-1.100) (-2.500,-1.000) (-2.595,-0.900) (-2.680,-0.800) (-2.755,-0.700) (-2.820,-0.600) (-2.875,-0.500) (-2.920,-0.400) (-2.955,-0.300) (-2.980,-0.200) (-2.995,-0.100) (-3.000,0.000) (-2.995,0.100) (-2.980,0.200) (-2.955,0.300) (-2.920,0.400) (-2.875,0.500) (-2.820,0.600) (-2.755,0.700) (-2.680,0.800) (-2.595,0.900) (-2.500,1.000) (-2.395,1.100) (-2.280,1.200) (-2.155,1.300) (-2.020,1.400) (-1.875,1.500) (-1.720,1.600) (-1.555,1.700) (-1.380,1.800) (-1.195,1.900) (-1.000,2.000) (-0.795,2.100) (-0.580,2.200) (-0.355,2.300) (-0.120,2.400) (0.125,2.500) (0.380,2.600) (0.645,2.700) (0.920,2.800) (1.205,2.900) (1.500,3.000) (1.805,3.100) (2.120,3.200) (2.445,3.300) (2.780,3.400) (3.125,3.500) (3.480,3.600) (3.845,3.700) (4.220,3.800) (4.605,3.900) (5.000,4.000)};
\end{scope}
\end{tikzpicture}
  \task \begin{tikzpicture}[x=0.27cm,y=0.27cm,baseline=(current bounding box.center)]
\draw[muted!22,very thin,xstep=1,ystep=1] (-5,-4) grid (5,4);
\draw[-latex,muted] (-5,0)--(5.35,0) node[right,font=\scriptsize]{$x$};
\draw[-latex,muted] (0,-4)--(0,4.35) node[above,font=\scriptsize]{$y$};
\node[font=\tiny,muted,below] at (-4,0) {-4};
\node[font=\tiny,muted,below] at (-2,0) {-2};
\node[font=\tiny,muted,below] at (2,0) {2};
\node[font=\tiny,muted,below] at (4,0) {4};
\node[font=\tiny,muted,left] at (0,-4) {-4};
\node[font=\tiny,muted,left] at (0,-2) {-2};
\node[font=\tiny,muted,left] at (0,2) {2};
\begin{scope}\clip (-5,-4) rectangle (5,4);
\draw[teal,line width=1.1pt] (-4.000,-3.000)--(0.000,3.000)--(-2.000,0.000)--(4.000,2.000);
\end{scope}
\end{tikzpicture}
\end{tasks}

\ex{2} ¿Las siguientes tablas definen una función de $x$?\begin{tasks}(2)
  \task $\begin{array}{|c|c|c|c|c|}\hline x&1&2&2&3\\\hline y&4&5&6&7\\\hline\end{array}$
  \task $\begin{array}{|c|c|c|c|c|}\hline x&1&2&3&4\\\hline y&5&5&5&6\\\hline\end{array}$
\end{tasks}

\ex{2} Explica por qué la recta $y=3$ es la gráfica de una función y la recta $x=3$ no lo es.

\ex{3} Dibuja una gráfica que pase por los puntos $(-2,\,1)$, $(0,\,3)$ y $(3,\,1)$ y que sea una función, y otra que pase por los mismos puntos y que \textbf{no} lo sea.

\ex{3} A cada número $x$ le hacemos corresponder el número $y$ que cumple $y^2=x$. ¿Es una función? ¿Y si la condición es $y^3=x$?


% =====================================================
\seccio{2}{Expresión de una función}

\ex{1} Dada la función $f(x)=3x-5$, calcula:\begin{tasks}(4)
  \task $f(0)$
  \task $f(2)$
  \task $f(-3)$
  \task $f\!\left(\tfrac12\right)$
\end{tasks}\recorda{Para calcular la imagen $f(a)$, sustituimos la $x$ por $a$ en la fórmula. Los números negativos, entre paréntesis.}

\ex{1} Si $g(x)=x^2-4x+3$, calcula $g(0)$, $g(1)$, $g(-2)$ y $g(5)$.

\ex{1} Completa las tablas de valores:\begin{tasks}(2)
  \task $y=2x-1$\\[3pt]$\begin{array}{|c|C{0.7cm}C{0.7cm}C{0.7cm}C{0.7cm}C{0.7cm}|}\hline x&-2&-1&0&1&2\\\hline y&\phantom{00}&\phantom{00}&\phantom{00}&\phantom{00}&\phantom{00}\\\hline\end{array}$
  \task $y=-x^2+3$\\[3pt]$\begin{array}{|c|C{0.7cm}C{0.7cm}C{0.7cm}C{0.7cm}C{0.7cm}|}\hline x&-2&-1&0&1&2\\\hline y&\phantom{00}&\phantom{00}&\phantom{00}&\phantom{00}&\phantom{00}\\\hline\end{array}$
\end{tasks}

\ex{2} Calcula las imágenes que se indican:\begin{tasks}(2)
  \task $h(x)=\dfrac{2x+1}{3}$: \ $h(4)$, $h(-2)$, $h(1)$
  \task $k(x)=\sqrt{x+5}$: \ $k(4)$, $k(-1)$, $k(11)$
\end{tasks}

\ex{2} Dada $f(x)=5x-3$, halla la antiimagen de:\begin{tasks}(3)
  \task $12$
  \task $-8$
  \task $0$
\end{tasks}\recorda{Para hallar la antiimagen de $k$, resolvemos la ecuación $f(x)=k$.}

\ex{2} Dada $f(x)=x^2-1$, ¿qué valores de $x$ tienen por imagen:\begin{tasks}(3)
  \task $8$
  \task $-1$
  \task $-5$
\end{tasks}

\ex{2} Escribe la fórmula de cada función:\begin{tasks}(1)
  \task El perímetro de un rectángulo de base $x$ y altura 4 cm.
  \task El precio de $x$ entradas de concierto a 8,50 € cada una.
  \task El área de un círculo de radio $x$.
  \task El tiempo que se tarda en hacer 300 km a una velocidad de $x$ km/h.
\end{tasks}

\ex{2} Halla la fórmula de las funciones que dan estas tablas:\begin{tasks}(2)
  \task $\begin{array}{|c|ccccc|}\hline x&0&1&2&3&4\\\hline y&3&7&11&15&19\\\hline\end{array}$
  \task $\begin{array}{|c|cccc|}\hline x&1&2&3&4\\\hline y&1&4&9&16\\\hline\end{array}$
\end{tasks}

\ex{3} Una tienda de fotografía cobra 2 € por el envío más 0,15 € por cada foto revelada. Escribe la fórmula del precio en función del número de fotos, calcula el precio de 40 fotos y cuántas fotos podemos revelar con 11 €.

\ex{3} La función $f(x)=ax+3$ cumple que $f(2)=11$. Halla $a$ y calcula $f(-1)$.


% =====================================================
\seccio{3}{Las funciones y sus gráficas}

\ex{1} Haz una tabla de valores de $f(x)=2x-3$ para $x=-1,\,0,\,1,\,2,\,3$, representa la función y halla los puntos de corte con los ejes.\recorda{Corte con el eje $Y$: hacemos $x=0$. Corte con el eje $X$: resolvemos $f(x)=0$.}

\ex{1} Haz una tabla de valores de $g(x)=x^2-4$ para $x$ entre $-3$ y $3$ y represéntala. ¿Dónde corta a los ejes?

\ex{2} Halla los puntos de corte con los ejes:\begin{tasks}(2)
  \task $f(x)=3x+6$
  \task $g(x)=x^2-5x+6$
  \task $h(x)=x^2+1$
  \task $k(x)=x^3-4x$
\end{tasks}

\ex{2} La gráfica muestra la distancia a casa de un ciclista durante una salida.\par\smallskip\begin{center}\begin{tikzpicture}[x=0.9cm,y=0.07cm,baseline=(current bounding box.center)]
\draw[muted!22,very thin,xstep=1,ystep=5] (0,0) grid (7,45);
\draw[-latex,muted] (0,0)--(7.35,0) node[right,font=\scriptsize]{$t\ (\mathrm{h})$};
\draw[-latex,muted] (0,0)--(0,46.75) node[above,font=\scriptsize]{$d\ (\mathrm{km})$};
\node[font=\tiny,muted,below] at (1,0) {1};
\node[font=\tiny,muted,below] at (2,0) {2};
\node[font=\tiny,muted,below] at (3,0) {3};
\node[font=\tiny,muted,below] at (4,0) {4};
\node[font=\tiny,muted,below] at (5,0) {5};
\node[font=\tiny,muted,below] at (6,0) {6};
\node[font=\tiny,muted,left] at (0,10) {10};
\node[font=\tiny,muted,left] at (0,20) {20};
\node[font=\tiny,muted,left] at (0,30) {30};
\node[font=\tiny,muted,left] at (0,40) {40};
\begin{scope}\clip (0,0) rectangle (7,45);
\draw[accent,line width=1.1pt] (0.0,0.0)--(2.0,30.0)--(3.0,30.0)--(4.0,40.0)--(6.0,0.0);
\end{scope}
\fill[accent] (0.0,0.0) circle (2.2pt);
\fill[accent] (6.0,0.0) circle (2.2pt);
\end{tikzpicture}\end{center}\begin{tasks}(1)
  \task ¿Cuántos kilómetros recorre en total?
  \task ¿Cuánto tiempo está parado?
  \task ¿A qué velocidad va en las dos primeras horas? ¿Y en la vuelta?
  \task ¿A qué distancia de casa está a las 5 h?
  \task ¿En qué momentos está a 15 km de casa?
\end{tasks}

\ex{2}\label{ex:temp} La gráfica muestra la temperatura de un día de invierno en una ciudad.\par\smallskip\begin{center}\begin{tikzpicture}[x=0.3cm,y=0.17cm,baseline=(current bounding box.center)]
\draw[muted!22,very thin,xstep=1,ystep=1] (0,-2) grid (25,18);
\draw[-latex,muted] (0,0)--(25.35,0) node[right,font=\scriptsize]{$h$};
\draw[-latex,muted] (0,-2)--(0,18.35) node[above,font=\scriptsize]{$T\ (^\circ\mathrm{C})$};
\node[font=\tiny,muted,below] at (2,0) {2};
\node[font=\tiny,muted,below] at (4,0) {4};
\node[font=\tiny,muted,below] at (6,0) {6};
\node[font=\tiny,muted,below] at (8,0) {8};
\node[font=\tiny,muted,below] at (10,0) {10};
\node[font=\tiny,muted,below] at (12,0) {12};
\node[font=\tiny,muted,below] at (14,0) {14};
\node[font=\tiny,muted,below] at (16,0) {16};
\node[font=\tiny,muted,below] at (18,0) {18};
\node[font=\tiny,muted,below] at (20,0) {20};
\node[font=\tiny,muted,below] at (22,0) {22};
\node[font=\tiny,muted,below] at (24,0) {24};
\node[font=\tiny,muted,left] at (0,-2) {-2};
\node[font=\tiny,muted,left] at (0,2) {2};
\node[font=\tiny,muted,left] at (0,4) {4};
\node[font=\tiny,muted,left] at (0,6) {6};
\node[font=\tiny,muted,left] at (0,8) {8};
\node[font=\tiny,muted,left] at (0,10) {10};
\node[font=\tiny,muted,left] at (0,12) {12};
\node[font=\tiny,muted,left] at (0,14) {14};
\node[font=\tiny,muted,left] at (0,16) {16};
\begin{scope}\clip (0,-2) rectangle (25,18);
\draw[accent,line width=1.1pt] (0.0,4.0)--(6.0,0.0)--(14.0,16.0)--(20.0,10.0)--(24.0,6.0);
\end{scope}
\fill[accent] (0.0,4.0) circle (2.2pt);
\fill[accent] (24.0,6.0) circle (2.2pt);
\end{tikzpicture}\end{center}\begin{tasks}(2)
  \task ¿Qué temperatura había a las 10 h?
  \task ¿A qué hora hizo más calor? ¿Y más frío?
  \task ¿Cuándo había 10 °C?
  \task ¿Cuánto subió la temperatura de 6 h a 14 h?
\end{tasks}

\ex{2} Asocia cada situación con su gráfica:\begin{tasks}(1)
  \task El precio de la fruta según los kilos que compramos.
  \task La altura de una pelota que lanzamos hacia arriba.
  \task La temperatura de un café caliente que se enfría.
  \task El agua de una bañera que llenamos, dejamos llena un rato y después vaciamos.
\end{tasks}\par\smallskip\begin{center}\textbf{I}\ \begin{tikzpicture}[x=0.3cm,y=0.3cm,baseline=(current bounding box.center)]
\draw[muted!22,very thin,xstep=1,ystep=1] (-1,-1) grid (6,6);
\draw[-latex,muted] (-1,0)--(6.35,0) node[right,font=\scriptsize]{$t$};
\draw[-latex,muted] (0,-1)--(0,6.35) node[above,font=\scriptsize]{$y$};
\begin{scope}\clip (-1,-1) rectangle (6,6);
\draw[accent,line width=1.1pt] plot[smooth] coordinates {(0.000,0.000) (0.031,0.124) (0.062,0.247) (0.094,0.368) (0.125,0.487) (0.156,0.605) (0.188,0.722) (0.219,0.837) (0.250,0.950) (0.281,1.062) (0.312,1.172) (0.344,1.280) (0.375,1.387) (0.406,1.493) (0.438,1.597) (0.469,1.699) (0.500,1.800) (0.531,1.899) (0.562,1.997) (0.594,2.093) (0.625,2.188) (0.656,2.280) (0.688,2.372) (0.719,2.462) (0.750,2.550) (0.781,2.637) (0.812,2.722) (0.844,2.805) (0.875,2.887) (0.906,2.968) (0.938,3.047) (0.969,3.124) (1.000,3.200) (1.031,3.274) (1.062,3.347) (1.094,3.418) (1.125,3.487) (1.156,3.555) (1.188,3.622) (1.219,3.687) (1.250,3.750) (1.281,3.812) (1.312,3.872) (1.344,3.930) (1.375,3.987) (1.406,4.043) (1.438,4.097) (1.469,4.149) (1.500,4.200) (1.531,4.249) (1.562,4.297) (1.594,4.343) (1.625,4.388) (1.656,4.430) (1.688,4.472) (1.719,4.512) (1.750,4.550) (1.781,4.587) (1.812,4.622) (1.844,4.655) (1.875,4.688) (1.906,4.718) (1.938,4.747) (1.969,4.774) (2.000,4.800) (2.031,4.824) (2.062,4.847) (2.094,4.868) (2.125,4.888) (2.156,4.905) (2.188,4.922) (2.219,4.937) (2.250,4.950) (2.281,4.962) (2.312,4.972) (2.344,4.980) (2.375,4.987) (2.406,4.993) (2.438,4.997) (2.469,4.999) (2.500,5.000) (2.531,4.999) (2.562,4.997) (2.594,4.993) (2.625,4.987) (2.656,4.980) (2.688,4.972) (2.719,4.962) (2.750,4.950) (2.781,4.937) (2.812,4.922) (2.844,4.905) (2.875,4.888) (2.906,4.868) (2.938,4.847) (2.969,4.824) (3.000,4.800) (3.031,4.774) (3.062,4.747) (3.094,4.718) (3.125,4.688) (3.156,4.655) (3.188,4.622) (3.219,4.587) (3.250,4.550) (3.281,4.512) (3.312,4.472) (3.344,4.430) (3.375,4.388) (3.406,4.343) (3.438,4.297) (3.469,4.249) (3.500,4.200) (3.531,4.149) (3.562,4.097) (3.594,4.043) (3.625,3.987) (3.656,3.930) (3.688,3.872) (3.719,3.812) (3.750,3.750) (3.781,3.687) (3.812,3.622) (3.844,3.555) (3.875,3.487) (3.906,3.418) (3.938,3.347) (3.969,3.274) (4.000,3.200) (4.031,3.124) (4.062,3.047) (4.094,2.968) (4.125,2.887) (4.156,2.805) (4.188,2.722) (4.219,2.637) (4.250,2.550) (4.281,2.462) (4.312,2.372) (4.344,2.280) (4.375,2.188) (4.406,2.093) (4.438,1.997) (4.469,1.899) (4.500,1.800) (4.531,1.699) (4.562,1.597) (4.594,1.493) (4.625,1.387) (4.656,1.280) (4.688,1.172) (4.719,1.062) (4.750,0.950) (4.781,0.837) (4.812,0.722) (4.844,0.605) (4.875,0.487) (4.906,0.368) (4.938,0.247) (4.969,0.124) (5.000,0.000)};
\end{scope}
\end{tikzpicture}\hfill\textbf{II}\ \begin{tikzpicture}[x=0.3cm,y=0.3cm,baseline=(current bounding box.center)]
\draw[muted!22,very thin,xstep=1,ystep=1] (-1,-1) grid (6,6);
\draw[-latex,muted] (-1,0)--(6.35,0) node[right,font=\scriptsize]{$t$};
\draw[-latex,muted] (0,-1)--(0,6.35) node[above,font=\scriptsize]{$y$};
\begin{scope}\clip (-1,-1) rectangle (6,6);
\draw[accent,line width=1.1pt] (0.0,0.0)--(2.0,5.0)--(4.0,5.0)--(5.0,0.0);
\end{scope}
\fill[accent] (0.0,0.0) circle (2.2pt);
\fill[accent] (5.0,0.0) circle (2.2pt);
\end{tikzpicture}\hfill\textbf{III}\ \begin{tikzpicture}[x=0.3cm,y=0.3cm,baseline=(current bounding box.center)]
\draw[muted!22,very thin,xstep=1,ystep=1] (-1,-1) grid (6,6);
\draw[-latex,muted] (-1,0)--(6.35,0) node[right,font=\scriptsize]{$x$};
\draw[-latex,muted] (0,-1)--(0,6.35) node[above,font=\scriptsize]{$y$};
\begin{scope}\clip (-1,-1) rectangle (6,6);
\draw[accent,line width=1.1pt] plot[smooth] coordinates {(0.000,0.000) (0.034,0.034) (0.069,0.069) (0.103,0.103) (0.138,0.138) (0.172,0.172) (0.206,0.206) (0.241,0.241) (0.275,0.275) (0.309,0.309) (0.344,0.344) (0.378,0.378) (0.412,0.412) (0.447,0.447) (0.481,0.481) (0.516,0.516) (0.550,0.550) (0.584,0.584) (0.619,0.619) (0.653,0.653) (0.688,0.688) (0.722,0.722) (0.756,0.756) (0.791,0.791) (0.825,0.825) (0.859,0.859) (0.894,0.894) (0.928,0.928) (0.963,0.963) (0.997,0.997) (1.031,1.031) (1.066,1.066) (1.100,1.100) (1.134,1.134) (1.169,1.169) (1.203,1.203) (1.238,1.238) (1.272,1.272) (1.306,1.306) (1.341,1.341) (1.375,1.375) (1.409,1.409) (1.444,1.444) (1.478,1.478) (1.512,1.512) (1.547,1.547) (1.581,1.581) (1.616,1.616) (1.650,1.650) (1.684,1.684) (1.719,1.719) (1.753,1.753) (1.788,1.788) (1.822,1.822) (1.856,1.856) (1.891,1.891) (1.925,1.925) (1.959,1.959) (1.994,1.994) (2.028,2.028) (2.062,2.062) (2.097,2.097) (2.131,2.131) (2.166,2.166) (2.200,2.200) (2.234,2.234) (2.269,2.269) (2.303,2.303) (2.337,2.337) (2.372,2.372) (2.406,2.406) (2.441,2.441) (2.475,2.475) (2.509,2.509) (2.544,2.544) (2.578,2.578) (2.612,2.612) (2.647,2.647) (2.681,2.681) (2.716,2.716) (2.750,2.750) (2.784,2.784) (2.819,2.819) (2.853,2.853) (2.888,2.888) (2.922,2.922) (2.956,2.956) (2.991,2.991) (3.025,3.025) (3.059,3.059) (3.094,3.094) (3.128,3.128) (3.163,3.163) (3.197,3.197) (3.231,3.231) (3.266,3.266) (3.300,3.300) (3.334,3.334) (3.369,3.369) (3.403,3.403) (3.438,3.438) (3.472,3.472) (3.506,3.506) (3.541,3.541) (3.575,3.575) (3.609,3.609) (3.644,3.644) (3.678,3.678) (3.712,3.712) (3.747,3.747) (3.781,3.781) (3.816,3.816) (3.850,3.850) (3.884,3.884) (3.919,3.919) (3.953,3.953) (3.987,3.987) (4.022,4.022) (4.056,4.056) (4.091,4.091) (4.125,4.125) (4.159,4.159) (4.194,4.194) (4.228,4.228) (4.263,4.263) (4.297,4.297) (4.331,4.331) (4.366,4.366) (4.400,4.400) (4.434,4.434) (4.469,4.469) (4.503,4.503) (4.537,4.537) (4.572,4.572) (4.606,4.606) (4.641,4.641) (4.675,4.675) (4.709,4.709) (4.744,4.744) (4.778,4.778) (4.812,4.812) (4.847,4.847) (4.881,4.881) (4.916,4.916) (4.950,4.950) (4.984,4.984) (5.019,5.019) (5.053,5.053) (5.088,5.088) (5.122,5.122) (5.156,5.156) (5.191,5.191) (5.225,5.225) (5.259,5.259) (5.294,5.294) (5.328,5.328) (5.362,5.362) (5.397,5.397) (5.431,5.431) (5.466,5.466) (5.500,5.500)};
\end{scope}
\end{tikzpicture}\hfill\textbf{IV}\ \begin{tikzpicture}[x=0.3cm,y=0.3cm,baseline=(current bounding box.center)]
\draw[muted!22,very thin,xstep=1,ystep=1] (-1,-1) grid (6,6);
\draw[-latex,muted] (-1,0)--(6.35,0) node[right,font=\scriptsize]{$t$};
\draw[-latex,muted] (0,-1)--(0,6.35) node[above,font=\scriptsize]{$y$};
\begin{scope}\clip (-1,-1) rectangle (6,6);
\draw[accent,line width=1.1pt] plot[smooth] coordinates {(0.000,5.000) (0.034,4.918) (0.069,4.838) (0.103,4.760) (0.138,4.683) (0.172,4.608) (0.206,4.534) (0.241,4.462) (0.275,4.392) (0.309,4.322) (0.344,4.255) (0.378,4.188) (0.412,4.123) (0.447,4.059) (0.481,3.997) (0.516,3.936) (0.550,3.876) (0.584,3.817) (0.619,3.759) (0.653,3.703) (0.688,3.648) (0.722,3.594) (0.756,3.541) (0.791,3.489) (0.825,3.438) (0.859,3.389) (0.894,3.340) (0.928,3.292) (0.963,3.245) (0.997,3.199) (1.031,3.154) (1.066,3.110) (1.100,3.067) (1.134,3.025) (1.169,2.984) (1.203,2.943) (1.238,2.904) (1.272,2.865) (1.306,2.827) (1.341,2.789) (1.375,2.753) (1.409,2.717) (1.444,2.682) (1.478,2.648) (1.512,2.614) (1.547,2.581) (1.581,2.549) (1.616,2.517) (1.650,2.486) (1.684,2.456) (1.719,2.426) (1.753,2.397) (1.788,2.369) (1.822,2.341) (1.856,2.313) (1.891,2.286) (1.925,2.260) (1.959,2.235) (1.994,2.209) (2.028,2.185) (2.062,2.160) (2.097,2.137) (2.131,2.114) (2.166,2.091) (2.200,2.069) (2.234,2.047) (2.269,2.025) (2.303,2.004) (2.337,1.984) (2.372,1.964) (2.406,1.944) (2.441,1.925) (2.475,1.906) (2.509,1.888) (2.544,1.869) (2.578,1.852) (2.612,1.834) (2.647,1.817) (2.681,1.801) (2.716,1.784) (2.750,1.768) (2.784,1.753) (2.819,1.737) (2.853,1.722) (2.888,1.707) (2.922,1.693) (2.956,1.679) (2.991,1.665) (3.025,1.651) (3.059,1.638) (3.094,1.625) (3.128,1.612) (3.163,1.600) (3.197,1.588) (3.231,1.576) (3.266,1.564) (3.300,1.552) (3.334,1.541) (3.369,1.530) (3.403,1.519) (3.438,1.509) (3.472,1.498) (3.506,1.488) (3.541,1.478) (3.575,1.468) (3.609,1.459) (3.644,1.449) (3.678,1.440) (3.712,1.431) (3.747,1.422) (3.781,1.414) (3.816,1.405) (3.850,1.397) (3.884,1.389) (3.919,1.381) (3.953,1.373) (3.987,1.366) (4.022,1.358) (4.056,1.351) (4.091,1.344) (4.125,1.337) (4.159,1.330) (4.194,1.323) (4.228,1.316) (4.263,1.310) (4.297,1.304) (4.331,1.297) (4.366,1.291) (4.400,1.285) (4.434,1.280) (4.469,1.274) (4.503,1.268) (4.537,1.263) (4.572,1.257) (4.606,1.252) (4.641,1.247) (4.675,1.242) (4.709,1.237) (4.744,1.232) (4.778,1.228) (4.812,1.223) (4.847,1.218) (4.881,1.214) (4.916,1.209) (4.950,1.205) (4.984,1.201) (5.019,1.197) (5.053,1.193) (5.088,1.189) (5.122,1.185) (5.156,1.181) (5.191,1.178) (5.225,1.174) (5.259,1.170) (5.294,1.167) (5.328,1.164) (5.362,1.160) (5.397,1.157) (5.431,1.154) (5.466,1.151) (5.500,1.148)};
\end{scope}
\end{tikzpicture}\end{center}

\ex{3} Haz una tabla de valores de $f(x)=\dfrac{12}{x}$ para $x=\pm1,\,\pm2,\,\pm3,\,\pm4,\,\pm6$ y represéntala. ¿Por qué no podemos unir el punto de $x=-1$ con el de $x=1$?

\ex{3} Inventa una historia que corresponda a esta gráfica de la distancia a casa de una persona.\par\smallskip\begin{center}\begin{tikzpicture}[x=0.8cm,y=0.5cm,baseline=(current bounding box.center)]
\draw[muted!22,very thin,xstep=1,ystep=1] (0,0) grid (6,4);
\draw[-latex,muted] (0,0)--(6.35,0) node[right,font=\scriptsize]{$t\ (\mathrm{h})$};
\draw[-latex,muted] (0,0)--(0,4.35) node[above,font=\scriptsize]{$d\ (\mathrm{km})$};
\node[font=\tiny,muted,below] at (1,0) {1};
\node[font=\tiny,muted,below] at (2,0) {2};
\node[font=\tiny,muted,below] at (3,0) {3};
\node[font=\tiny,muted,below] at (4,0) {4};
\node[font=\tiny,muted,below] at (5,0) {5};
\node[font=\tiny,muted,left] at (0,1) {1};
\node[font=\tiny,muted,left] at (0,2) {2};
\node[font=\tiny,muted,left] at (0,3) {3};
\begin{scope}\clip (0,0) rectangle (6,4);
\draw[accent,line width=1.1pt] (0.0,0.0)--(1.0,2.0)--(3.0,2.0)--(3.5,0.0)--(4.0,0.0)--(5.0,3.0);
\end{scope}
\fill[accent] (0.0,0.0) circle (2.2pt);
\fill[accent] (5.0,3.0) circle (2.2pt);
\end{tikzpicture}\end{center}


% =====================================================
\seccio{4}{Dominio y recorrido}

\ex{1} Escribe el dominio y el recorrido de estas funciones:\begin{tasks}(3)
  \task \begin{tikzpicture}[x=0.27cm,y=0.27cm,baseline=(current bounding box.center)]
\draw[muted!22,very thin,xstep=1,ystep=1] (-5,-3) grid (5,4);
\draw[-latex,muted] (-5,0)--(5.35,0) node[right,font=\scriptsize]{$x$};
\draw[-latex,muted] (0,-3)--(0,4.35) node[above,font=\scriptsize]{$y$};
\node[font=\tiny,muted,below] at (-4,0) {-4};
\node[font=\tiny,muted,below] at (-2,0) {-2};
\node[font=\tiny,muted,below] at (2,0) {2};
\node[font=\tiny,muted,below] at (4,0) {4};
\node[font=\tiny,muted,left] at (0,-2) {-2};
\node[font=\tiny,muted,left] at (0,2) {2};
\begin{scope}\clip (-5,-3) rectangle (5,4);
\draw[accent,line width=1.1pt] (-4.0,3.0)--(-1.0,-2.0)--(2.0,1.0)--(4.0,1.0);
\end{scope}
\fill[accent] (-4.0,3.0) circle (2.2pt);
\fill[accent] (4.0,1.0) circle (2.2pt);
\end{tikzpicture}
  \task \begin{tikzpicture}[x=0.27cm,y=0.27cm,baseline=(current bounding box.center)]
\draw[muted!22,very thin,xstep=1,ystep=1] (-4,-3) grid (5,4);
\draw[-latex,muted] (-4,0)--(5.35,0) node[right,font=\scriptsize]{$x$};
\draw[-latex,muted] (0,-3)--(0,4.35) node[above,font=\scriptsize]{$y$};
\node[font=\tiny,muted,below] at (-4,0) {-4};
\node[font=\tiny,muted,below] at (-2,0) {-2};
\node[font=\tiny,muted,below] at (2,0) {2};
\node[font=\tiny,muted,below] at (4,0) {4};
\node[font=\tiny,muted,left] at (0,-2) {-2};
\node[font=\tiny,muted,left] at (0,2) {2};
\begin{scope}\clip (-4,-3) rectangle (5,4);
\draw[accent,line width=1.1pt] (-3.0,-1.0)--(1.0,3.0);
\draw[accent,line width=1.1pt] (1.0,0.0)--(4.0,-2.0);
\end{scope}
\filldraw[fill=white,draw=accent,line width=0.9pt] (-3.0,-1.0) circle (2.2pt);
\fill[accent] (1.0,3.0) circle (2.2pt);
\filldraw[fill=white,draw=accent,line width=0.9pt] (1.0,0.0) circle (2.2pt);
\filldraw[fill=white,draw=accent,line width=0.9pt] (4.0,-2.0) circle (2.2pt);
\end{tikzpicture}
  \task \begin{tikzpicture}[x=0.27cm,y=0.27cm,baseline=(current bounding box.center)]
\draw[muted!22,very thin,xstep=1,ystep=1] (-4,-2) grid (4,3);
\draw[-latex,muted] (-4,0)--(4.35,0) node[right,font=\scriptsize]{$x$};
\draw[-latex,muted] (0,-2)--(0,3.35) node[above,font=\scriptsize]{$y$};
\node[font=\tiny,muted,below] at (-4,0) {-4};
\node[font=\tiny,muted,below] at (-2,0) {-2};
\node[font=\tiny,muted,below] at (2,0) {2};
\node[font=\tiny,muted,left] at (0,-2) {-2};
\node[font=\tiny,muted,left] at (0,2) {2};
\begin{scope}\clip (-4,-2) rectangle (4,3);
\draw[accent,line width=1.1pt] (-2.0,2.0)--(0.0,-1.0)--(2.0,2.0);
\draw[accent,line width=1.1pt,-latex] (-2.0,2.0)--(-2.467,2.700);
\draw[accent,line width=1.1pt,-latex] (2.0,2.0)--(2.467,2.700);
\end{scope}
\end{tikzpicture}
\end{tasks}\recorda{Dominio: valores de $x$ que tienen imagen (proyección sobre el eje $X$). Recorrido: valores que toma la $y$ (proyección sobre el eje $Y$). Punto lleno: $[\ ]$; punto vacío: $(\ )$; flecha: $\pm\infty$.}

\ex{1} Calcula el dominio de:\begin{tasks}(4)
  \task $f(x)=3x^2-2x+1$
  \task $g(x)=\dfrac{4}{x-1}$
  \task $h(x)=\sqrt{x-2}$
  \task $k(x)=\dfrac{x+3}{x+5}$
\end{tasks}

\ex{2} Escribe el dominio y el recorrido de estas funciones:\begin{tasks}(3)
  \task \begin{tikzpicture}[x=0.25cm,y=0.25cm,baseline=(current bounding box.center)]
\draw[muted!22,very thin,xstep=1,ystep=1] (-6,-3) grid (6,5);
\draw[-latex,muted] (-6,0)--(6.35,0) node[right,font=\scriptsize]{$x$};
\draw[-latex,muted] (0,-3)--(0,5.35) node[above,font=\scriptsize]{$y$};
\node[font=\tiny,muted,below] at (-6,0) {-6};
\node[font=\tiny,muted,below] at (-4,0) {-4};
\node[font=\tiny,muted,below] at (-2,0) {-2};
\node[font=\tiny,muted,below] at (2,0) {2};
\node[font=\tiny,muted,below] at (4,0) {4};
\node[font=\tiny,muted,left] at (0,-2) {-2};
\node[font=\tiny,muted,left] at (0,2) {2};
\node[font=\tiny,muted,left] at (0,4) {4};
\begin{scope}\clip (-6,-3) rectangle (6,5);
\draw[accent,line width=1.1pt] (-5.0,1.0)--(-2.0,4.0);
\draw[accent,line width=1.1pt] (1.0,-2.0)--(4.0,0.0);
\draw[accent,line width=1.1pt,-latex] (4.0,0.0)--(5.800,1.200);
\end{scope}
\fill[accent] (-5.0,1.0) circle (2.2pt);
\fill[accent] (-2.0,4.0) circle (2.2pt);
\fill[accent] (1.0,-2.0) circle (2.2pt);
\end{tikzpicture}
  \task \begin{tikzpicture}[x=0.27cm,y=0.27cm,baseline=(current bounding box.center)]
\draw[muted!22,very thin,xstep=1,ystep=1] (-5,-4) grid (6,3);
\draw[-latex,muted] (-5,0)--(6.35,0) node[right,font=\scriptsize]{$x$};
\draw[-latex,muted] (0,-4)--(0,3.35) node[above,font=\scriptsize]{$y$};
\node[font=\tiny,muted,below] at (-4,0) {-4};
\node[font=\tiny,muted,below] at (-2,0) {-2};
\node[font=\tiny,muted,below] at (2,0) {2};
\node[font=\tiny,muted,below] at (4,0) {4};
\node[font=\tiny,muted,left] at (0,-4) {-4};
\node[font=\tiny,muted,left] at (0,-2) {-2};
\node[font=\tiny,muted,left] at (0,2) {2};
\begin{scope}\clip (-5,-4) rectangle (6,3);
\draw[accent,line width=1.1pt] (-4.0,-3.0)--(-1.0,0.0);
\draw[accent,line width=1.1pt] (-1.0,2.0)--(3.0,2.0);
\end{scope}
\fill[accent] (-4.0,-3.0) circle (2.2pt);
\filldraw[fill=white,draw=accent,line width=0.9pt] (-1.0,0.0) circle (2.2pt);
\fill[accent] (-1.0,2.0) circle (2.2pt);
\filldraw[fill=white,draw=accent,line width=0.9pt] (3.0,2.0) circle (2.2pt);
\fill[accent] (5.0,-1.0) circle (2.2pt);
\end{tikzpicture}
  \task \begin{tikzpicture}[x=0.27cm,y=0.27cm,baseline=(current bounding box.center)]
\draw[muted!22,very thin,xstep=1,ystep=1] (-5,-3) grid (5,3);
\draw[-latex,muted] (-5,0)--(5.35,0) node[right,font=\scriptsize]{$x$};
\draw[-latex,muted] (0,-3)--(0,3.35) node[above,font=\scriptsize]{$y$};
\node[font=\tiny,muted,below] at (-4,0) {-4};
\node[font=\tiny,muted,below] at (-2,0) {-2};
\node[font=\tiny,muted,below] at (2,0) {2};
\node[font=\tiny,muted,below] at (4,0) {4};
\node[font=\tiny,muted,left] at (0,-2) {-2};
\node[font=\tiny,muted,left] at (0,2) {2};
\begin{scope}\clip (-5,-3) rectangle (5,3);
\draw[accent,line width=1.1pt] (-3.0,-1.0)--(0.0,2.0);
\draw[accent,line width=1.1pt,-latex] (-3.0,-1.0)--(-4.800,-2.800);
\draw[accent,line width=1.1pt] (0.0,-2.0)--(3.0,1.0);
\draw[accent,line width=1.1pt,-latex] (3.0,1.0)--(4.800,2.800);
\end{scope}
\filldraw[fill=white,draw=accent,line width=0.9pt] (0.0,2.0) circle (2.2pt);
\filldraw[fill=white,draw=accent,line width=0.9pt] (0.0,-2.0) circle (2.2pt);
\end{tikzpicture}
\end{tasks}

\ex{2} Calcula el dominio de:\begin{tasks}(4)
  \task $f(x)=\dfrac{1}{x^2-4}$
  \task $g(x)=\sqrt{6-2x}$
  \task $h(x)=\dfrac{x}{x^2+1}$
  \task $k(x)=\dfrac{5}{3x-9}$
\end{tasks}

\ex{2} ¿Qué dominio tienen estas funciones en su contexto?\begin{tasks}(1)
  \task El precio de $x$ bolígrafos a 0,80 € la unidad.
  \task El área de un cuadrado de lado $x$: $A(x)=x^2$.
  \task La altura de una pelota, $h(t)=20t-5t^2$, mientras está en el aire ($t$ en segundos).
\end{tasks}

\ex{3} Dibuja la gráfica de una función que tenga dominio $[-3,\,4]$ y recorrido $[-2,\,5]$.

\ex{3} Calcula el dominio de $f(x)=\dfrac{\sqrt{x+1}}{x-3}$.


% =====================================================
\seccio{5}{Crecimiento y decrecimiento}

\ex{1}\label{ex:mono1} Indica los intervalos donde estas funciones son crecientes, decrecientes o constantes:\begin{tasks}(3)
  \task \begin{tikzpicture}[x=0.27cm,y=0.27cm,baseline=(current bounding box.center)]
\draw[muted!22,very thin,xstep=1,ystep=1] (-6,-3) grid (5,4);
\draw[-latex,muted] (-6,0)--(5.35,0) node[right,font=\scriptsize]{$x$};
\draw[-latex,muted] (0,-3)--(0,4.35) node[above,font=\scriptsize]{$y$};
\node[font=\tiny,muted,below] at (-6,0) {-6};
\node[font=\tiny,muted,below] at (-4,0) {-4};
\node[font=\tiny,muted,below] at (-2,0) {-2};
\node[font=\tiny,muted,below] at (2,0) {2};
\node[font=\tiny,muted,below] at (4,0) {4};
\node[font=\tiny,muted,left] at (0,-2) {-2};
\node[font=\tiny,muted,left] at (0,2) {2};
\begin{scope}\clip (-6,-3) rectangle (5,4);
\draw[accent,line width=1.1pt] (-5.0,-2.0)--(-2.0,3.0)--(1.0,-1.0)--(4.0,2.0);
\end{scope}
\fill[accent] (-5.0,-2.0) circle (2.2pt);
\fill[accent] (4.0,2.0) circle (2.2pt);
\end{tikzpicture}
  \task \begin{tikzpicture}[x=0.27cm,y=0.27cm,baseline=(current bounding box.center)]
\draw[muted!22,very thin,xstep=1,ystep=1] (-5,-4) grid (5,5);
\draw[-latex,muted] (-5,0)--(5.35,0) node[right,font=\scriptsize]{$x$};
\draw[-latex,muted] (0,-4)--(0,5.35) node[above,font=\scriptsize]{$y$};
\node[font=\tiny,muted,below] at (-4,0) {-4};
\node[font=\tiny,muted,below] at (-2,0) {-2};
\node[font=\tiny,muted,below] at (2,0) {2};
\node[font=\tiny,muted,below] at (4,0) {4};
\node[font=\tiny,muted,left] at (0,-4) {-4};
\node[font=\tiny,muted,left] at (0,-2) {-2};
\node[font=\tiny,muted,left] at (0,2) {2};
\node[font=\tiny,muted,left] at (0,4) {4};
\begin{scope}\clip (-5,-4) rectangle (5,5);
\draw[accent,line width=1.1pt] (-4.0,4.0)--(-1.0,1.0)--(2.0,1.0)--(4.0,-3.0);
\end{scope}
\fill[accent] (-4.0,4.0) circle (2.2pt);
\fill[accent] (4.0,-3.0) circle (2.2pt);
\end{tikzpicture}
  \task \begin{tikzpicture}[x=0.27cm,y=0.27cm,baseline=(current bounding box.center)]
\draw[muted!22,very thin,xstep=1,ystep=1] (-3,-4) grid (5,2);
\draw[-latex,muted] (-3,0)--(5.35,0) node[right,font=\scriptsize]{$x$};
\draw[-latex,muted] (0,-4)--(0,2.35) node[above,font=\scriptsize]{$y$};
\node[font=\tiny,muted,below] at (-2,0) {-2};
\node[font=\tiny,muted,below] at (2,0) {2};
\node[font=\tiny,muted,below] at (4,0) {4};
\node[font=\tiny,muted,left] at (0,-4) {-4};
\node[font=\tiny,muted,left] at (0,-2) {-2};
\begin{scope}\clip (-3,-4) rectangle (5,2);
\draw[accent,line width=1.1pt] (-1.0,-3.0)--(1.0,1.0)--(3.0,-1.0);
\draw[accent,line width=1.1pt,-latex] (-1.0,-3.0)--(-1.300,-3.600);
\draw[accent,line width=1.1pt,-latex] (3.0,-1.0)--(4.800,-2.800);
\end{scope}
\end{tikzpicture}
\end{tasks}\recorda{Leemos la gráfica de izquierda a derecha. Los intervalos se dan con valores de la $x$ y son abiertos.}

\ex{1} Sin dibujarlas, di si estas funciones son crecientes, decrecientes o constantes:\begin{tasks}(4)
  \task $y=4x-1$
  \task $y=-3x+2$
  \task $y=5$
  \task $y=\tfrac{x}{2}$
\end{tasks}

\ex{2} Estudia el crecimiento y el decrecimiento de estas funciones:\begin{tasks}(2)
  \task \begin{tikzpicture}[x=0.3cm,y=0.3cm,baseline=(current bounding box.center)]
\draw[muted!22,very thin,xstep=1,ystep=1] (-5,-3) grid (4,3);
\draw[-latex,muted] (-5,0)--(4.35,0) node[right,font=\scriptsize]{$x$};
\draw[-latex,muted] (0,-3)--(0,3.35) node[above,font=\scriptsize]{$y$};
\node[font=\tiny,muted,below] at (-4,0) {-4};
\node[font=\tiny,muted,below] at (-2,0) {-2};
\node[font=\tiny,muted,below] at (2,0) {2};
\node[font=\tiny,muted,left] at (0,-2) {-2};
\node[font=\tiny,muted,left] at (0,2) {2};
\begin{scope}\clip (-5,-3) rectangle (4,3);
\draw[accent,line width=1.1pt] (-3.0,2.0)--(-1.0,-2.0)--(1.0,-2.0)--(3.0,2.0);
\draw[accent,line width=1.1pt,-latex] (-3.0,2.0)--(-3.300,2.600);
\end{scope}
\fill[accent] (3.0,2.0) circle (2.2pt);
\end{tikzpicture}
  \task \begin{tikzpicture}[x=0.3cm,y=0.3cm,baseline=(current bounding box.center)]
\draw[muted!22,very thin,xstep=1,ystep=1] (-5,-3) grid (6,3);
\draw[-latex,muted] (-5,0)--(6.35,0) node[right,font=\scriptsize]{$x$};
\draw[-latex,muted] (0,-3)--(0,3.35) node[above,font=\scriptsize]{$y$};
\node[font=\tiny,muted,below] at (-4,0) {-4};
\node[font=\tiny,muted,below] at (-2,0) {-2};
\node[font=\tiny,muted,below] at (2,0) {2};
\node[font=\tiny,muted,below] at (4,0) {4};
\node[font=\tiny,muted,left] at (0,-2) {-2};
\node[font=\tiny,muted,left] at (0,2) {2};
\begin{scope}\clip (-5,-3) rectangle (6,3);
\draw[accent,line width=1.1pt] (-4.0,-1.0)--(-1.0,2.0);
\draw[accent,line width=1.1pt] (-1.0,-2.0)--(2.0,1.0)--(5.0,-1.0);
\end{scope}
\fill[accent] (-4.0,-1.0) circle (2.2pt);
\filldraw[fill=white,draw=accent,line width=0.9pt] (-1.0,2.0) circle (2.2pt);
\fill[accent] (-1.0,-2.0) circle (2.2pt);
\fill[accent] (5.0,-1.0) circle (2.2pt);
\end{tikzpicture}
\end{tasks}

\ex{2} Con la gráfica de la temperatura del ejercicio \ref{ex:temp}, escribe los intervalos de tiempo en los que la temperatura sube y en los que baja.

\ex{3} Pablo dice que la función del apartado a) del ejercicio \ref{ex:mono1} «crece en $(-2,\,3)$». ¿Qué error ha cometido?

\ex{3} Dibuja la gráfica de una función de dominio $\mathbb{R}$ que sea decreciente en $(-\infty,\,-1)$, constante en $(-1,\,2)$ y creciente en $(2,\,+\infty)$.


% =====================================================
\seccio{6}{Máximos y mínimos}

\ex{1} Halla los máximos y mínimos relativos de estas funciones:\begin{tasks}(2)
  \task \begin{tikzpicture}[x=0.3cm,y=0.3cm,baseline=(current bounding box.center)]
\draw[muted!22,very thin,xstep=1,ystep=1] (-6,-3) grid (6,5);
\draw[-latex,muted] (-6,0)--(6.35,0) node[right,font=\scriptsize]{$x$};
\draw[-latex,muted] (0,-3)--(0,5.35) node[above,font=\scriptsize]{$y$};
\node[font=\tiny,muted,below] at (-6,0) {-6};
\node[font=\tiny,muted,below] at (-4,0) {-4};
\node[font=\tiny,muted,below] at (-2,0) {-2};
\node[font=\tiny,muted,below] at (2,0) {2};
\node[font=\tiny,muted,below] at (4,0) {4};
\node[font=\tiny,muted,left] at (0,-2) {-2};
\node[font=\tiny,muted,left] at (0,2) {2};
\node[font=\tiny,muted,left] at (0,4) {4};
\begin{scope}\clip (-6,-3) rectangle (6,5);
\draw[accent,line width=1.1pt] (-5.0,1.0)--(-3.0,4.0)--(0.0,-2.0)--(3.0,3.0)--(5.0,0.0);
\end{scope}
\fill[accent] (-5.0,1.0) circle (2.2pt);
\fill[accent] (5.0,0.0) circle (2.2pt);
\end{tikzpicture}
  \task \begin{tikzpicture}[x=0.3cm,y=0.3cm,baseline=(current bounding box.center)]
\draw[muted!22,very thin,xstep=1,ystep=1] (-5,-4) grid (4,4);
\draw[-latex,muted] (-5,0)--(4.35,0) node[right,font=\scriptsize]{$x$};
\draw[-latex,muted] (0,-4)--(0,4.35) node[above,font=\scriptsize]{$y$};
\node[font=\tiny,muted,below] at (-4,0) {-4};
\node[font=\tiny,muted,below] at (-2,0) {-2};
\node[font=\tiny,muted,below] at (2,0) {2};
\node[font=\tiny,muted,left] at (0,-4) {-4};
\node[font=\tiny,muted,left] at (0,-2) {-2};
\node[font=\tiny,muted,left] at (0,2) {2};
\begin{scope}\clip (-5,-4) rectangle (4,4);
\draw[accent,line width=1.1pt] (-3.0,-1.0)--(-1.0,3.0)--(2.0,-3.0);
\draw[accent,line width=1.1pt,-latex] (-3.0,-1.0)--(-4.300,-3.600);
\draw[accent,line width=1.1pt,-latex] (2.0,-3.0)--(2.300,-3.600);
\end{scope}
\end{tikzpicture}
\end{tasks}\recorda{Máximo relativo: la función pasa de creciente a decreciente. Mínimo relativo: pasa de decreciente a creciente. Se escriben como puntos $(x,\,y)$.}

\ex{2} Halla los máximos y mínimos, relativos y absolutos, de estas funciones:\begin{tasks}(2)
  \task \begin{tikzpicture}[x=0.3cm,y=0.3cm,baseline=(current bounding box.center)]
\draw[muted!22,very thin,xstep=1,ystep=1] (-5,-3) grid (6,4);
\draw[-latex,muted] (-5,0)--(6.35,0) node[right,font=\scriptsize]{$x$};
\draw[-latex,muted] (0,-3)--(0,4.35) node[above,font=\scriptsize]{$y$};
\node[font=\tiny,muted,below] at (-4,0) {-4};
\node[font=\tiny,muted,below] at (-2,0) {-2};
\node[font=\tiny,muted,below] at (2,0) {2};
\node[font=\tiny,muted,below] at (4,0) {4};
\node[font=\tiny,muted,left] at (0,-2) {-2};
\node[font=\tiny,muted,left] at (0,2) {2};
\begin{scope}\clip (-5,-3) rectangle (6,4);
\draw[accent,line width=1.1pt] (-4.0,-2.0)--(-1.0,3.0)--(1.0,1.0)--(3.0,3.0)--(5.0,-1.0);
\end{scope}
\fill[accent] (-4.0,-2.0) circle (2.2pt);
\filldraw[fill=white,draw=accent,line width=0.9pt] (5.0,-1.0) circle (2.2pt);
\end{tikzpicture}
  \task \begin{tikzpicture}[x=0.3cm,y=0.3cm,baseline=(current bounding box.center)]
\draw[muted!22,very thin,xstep=1,ystep=1] (-4,-2) grid (4,5);
\draw[-latex,muted] (-4,0)--(4.35,0) node[right,font=\scriptsize]{$x$};
\draw[-latex,muted] (0,-2)--(0,5.35) node[above,font=\scriptsize]{$y$};
\node[font=\tiny,muted,below] at (-4,0) {-4};
\node[font=\tiny,muted,below] at (-2,0) {-2};
\node[font=\tiny,muted,below] at (2,0) {2};
\node[font=\tiny,muted,left] at (0,-2) {-2};
\node[font=\tiny,muted,left] at (0,2) {2};
\node[font=\tiny,muted,left] at (0,4) {4};
\begin{scope}\clip (-4,-2) rectangle (4,5);
\draw[accent,line width=1.1pt] (-3.0,4.0)--(0.0,-1.0)--(3.0,4.0);
\end{scope}
\filldraw[fill=white,draw=accent,line width=0.9pt] (-3.0,4.0) circle (2.2pt);
\filldraw[fill=white,draw=accent,line width=0.9pt] (3.0,4.0) circle (2.2pt);
\end{tikzpicture}
\end{tasks}

\ex{2} Haz una tabla de valores de $f(x)=x^2-6x+5$ para $x=0,\,1,\ldots,\,6$ y represéntala. ¿Dónde tiene el mínimo? ¿Cuáles son los puntos de corte con los ejes?

\ex{2} La altura en metros de un cohete de agua es $h(t)=40t-5t^2$, con $t$ en segundos. Haz una tabla para $t=0,\,2,\,4,\,6,\,8$. ¿Qué altura máxima alcanza y cuándo? ¿Cuándo vuelve al suelo?

\ex{3} Dibuja una función de dominio $\mathbb{R}$ con un máximo relativo en $(-2,\,3)$ y un mínimo relativo en $(1,\,-1)$, que no tenga ni máximo ni mínimo absolutos.

\ex{3} Razona si son verdaderas o falsas:\begin{tasks}(1)
  \task Un máximo absoluto siempre es también un máximo relativo.
  \task Si una función continua tiene dos máximos relativos, entre ellos tiene que haber al menos un mínimo relativo.
  \task Si una función es creciente en todo el dominio, no tiene máximos ni mínimos relativos.
\end{tasks}


% =====================================================
\seccio{7}{Continuidad}

\ex{1} Di si estas funciones son continuas. Si no lo son, indica dónde son discontinuas y por qué:\begin{tasks}(3)
  \task \begin{tikzpicture}[x=0.27cm,y=0.27cm,baseline=(current bounding box.center)]
\draw[muted!22,very thin,xstep=1,ystep=1] (-5,-3) grid (5,3);
\draw[-latex,muted] (-5,0)--(5.35,0) node[right,font=\scriptsize]{$x$};
\draw[-latex,muted] (0,-3)--(0,3.35) node[above,font=\scriptsize]{$y$};
\node[font=\tiny,muted,below] at (-4,0) {-4};
\node[font=\tiny,muted,below] at (-2,0) {-2};
\node[font=\tiny,muted,below] at (2,0) {2};
\node[font=\tiny,muted,below] at (4,0) {4};
\node[font=\tiny,muted,left] at (0,-2) {-2};
\node[font=\tiny,muted,left] at (0,2) {2};
\begin{scope}\clip (-5,-3) rectangle (5,3);
\draw[accent,line width=1.1pt] (-4.0,-2.0)--(-1.0,1.0)--(2.0,-1.0)--(4.0,2.0);
\end{scope}
\fill[accent] (-4.0,-2.0) circle (2.2pt);
\fill[accent] (4.0,2.0) circle (2.2pt);
\end{tikzpicture}
  \task \begin{tikzpicture}[x=0.27cm,y=0.27cm,baseline=(current bounding box.center)]
\draw[muted!22,very thin,xstep=1,ystep=1] (-5,-2) grid (5,4);
\draw[-latex,muted] (-5,0)--(5.35,0) node[right,font=\scriptsize]{$x$};
\draw[-latex,muted] (0,-2)--(0,4.35) node[above,font=\scriptsize]{$y$};
\node[font=\tiny,muted,below] at (-4,0) {-4};
\node[font=\tiny,muted,below] at (-2,0) {-2};
\node[font=\tiny,muted,below] at (2,0) {2};
\node[font=\tiny,muted,below] at (4,0) {4};
\node[font=\tiny,muted,left] at (0,-2) {-2};
\node[font=\tiny,muted,left] at (0,2) {2};
\begin{scope}\clip (-5,-2) rectangle (5,4);
\draw[accent,line width=1.1pt] (-4.0,2.0)--(0.0,2.0);
\draw[accent,line width=1.1pt] (0.0,-1.0)--(4.0,3.0);
\end{scope}
\fill[accent] (-4.0,2.0) circle (2.2pt);
\filldraw[fill=white,draw=accent,line width=0.9pt] (0.0,2.0) circle (2.2pt);
\fill[accent] (0.0,-1.0) circle (2.2pt);
\fill[accent] (4.0,3.0) circle (2.2pt);
\end{tikzpicture}
  \task \begin{tikzpicture}[x=0.27cm,y=0.27cm,baseline=(current bounding box.center)]
\draw[muted!22,very thin,xstep=1,ystep=1] (-5,-2) grid (5,5);
\draw[-latex,muted] (-5,0)--(5.35,0) node[right,font=\scriptsize]{$x$};
\draw[-latex,muted] (0,-2)--(0,5.35) node[above,font=\scriptsize]{$y$};
\node[font=\tiny,muted,below] at (-4,0) {-4};
\node[font=\tiny,muted,below] at (-2,0) {-2};
\node[font=\tiny,muted,below] at (2,0) {2};
\node[font=\tiny,muted,below] at (4,0) {4};
\node[font=\tiny,muted,left] at (0,-2) {-2};
\node[font=\tiny,muted,left] at (0,2) {2};
\node[font=\tiny,muted,left] at (0,4) {4};
\begin{scope}\clip (-5,-2) rectangle (5,5);
\draw[accent,line width=1.1pt] (-4.0,-1.0)--(1.0,4.0);
\draw[accent,line width=1.1pt] (1.0,4.0)--(4.0,1.0);
\end{scope}
\fill[accent] (-4.0,-1.0) circle (2.2pt);
\filldraw[fill=white,draw=accent,line width=0.9pt] (1.0,4.0) circle (2.2pt);
\filldraw[fill=white,draw=accent,line width=0.9pt] (1.0,4.0) circle (2.2pt);
\fill[accent] (4.0,1.0) circle (2.2pt);
\fill[accent] (1.0,1.0) circle (2.2pt);
\end{tikzpicture}
\end{tasks}\recorda{Una función es continua si se puede dibujar sin levantar el lápiz del papel.}

\ex{1} Un aparcamiento cobra 1,50 € por cada hora o fracción, hasta un máximo de 4 horas. Representa la función precio y calcula cuánto se paga por 2 h 30 min. ¿Es continua?

\ex{2} Dada $f(x)=\begin{cases}x+3 & \text{si } x<0\\ 3-x^2 & \text{si } x\ge0\end{cases}$, calcula $f(-2)$, $f(0)$ y $f(2)$. ¿Es continua en $x=0$?

\ex{2} Dada $g(x)=\begin{cases}2x & \text{si } x\le1\\ x+3 & \text{si } x>1\end{cases}$, calcula $g(1)$ y di si es continua en $x=1$. Dibújala.

\ex{2} Una empresa de mensajería cobra 3 € por paquetes de hasta 1 kg, 5 € por paquetes de más de 1 kg hasta 3 kg y 8 € por paquetes de más de 3 kg hasta 5 kg. Dibuja la gráfica, escribe el dominio y el recorrido e indica dónde es discontinua.

\ex{3} Halla el valor de $a$ para que la función $f(x)=\begin{cases}x^2+a & \text{si } x<2\\ 3x-1 & \text{si } x\ge2\end{cases}$ sea continua.

\ex{3} ¿Dónde son discontinuas estas funciones? ¿Por qué?\begin{tasks}(2)
  \task $f(x)=\dfrac{1}{x-2}$
  \task $g(x)=\dfrac{3}{x^2-9}$
\end{tasks}


% =====================================================
\seccio{8}{Problemas con funciones}

\ex{1} Un taxi cobra 3,50 € por la bajada de bandera y 1,10 € por kilómetro. Escribe la fórmula del precio, calcula cuánto cuesta un trayecto de 12 km y cuántos kilómetros hemos hecho si hemos pagado 25,50 €.\recorda{Identifica las variables (qué es la $x$ y qué es la $y$, con las unidades), escribe la fórmula o dibuja la gráfica e interpreta cada resultado en el contexto del problema.}

\ex{1} Una tarifa de móvil A cuesta 10 € al mes más 0,05 € por minuto, y una tarifa B cuesta 0,15 € por minuto. Escribe las dos fórmulas. ¿Cuál le conviene a quien habla 60 minutos al mes? ¿Y 150 minutos? ¿Para cuántos minutos cuestan lo mismo?

\ex{2} Haz el estudio completo de esta función: dominio, recorrido, puntos de corte, crecimiento, máximos y mínimos y continuidad.\par\smallskip\begin{center}\begin{tikzpicture}[x=0.4cm,y=0.4cm,baseline=(current bounding box.center)]
\draw[muted!22,very thin,xstep=1,ystep=1] (-6,-5) grid (6,4);
\draw[-latex,muted] (-6,0)--(6.35,0) node[right,font=\scriptsize]{$x$};
\draw[-latex,muted] (0,-5)--(0,4.35) node[above,font=\scriptsize]{$y$};
\node[font=\tiny,muted,below] at (-6,0) {-6};
\node[font=\tiny,muted,below] at (-4,0) {-4};
\node[font=\tiny,muted,below] at (-2,0) {-2};
\node[font=\tiny,muted,below] at (2,0) {2};
\node[font=\tiny,muted,below] at (4,0) {4};
\node[font=\tiny,muted,left] at (0,-4) {-4};
\node[font=\tiny,muted,left] at (0,-2) {-2};
\node[font=\tiny,muted,left] at (0,2) {2};
\begin{scope}\clip (-6,-5) rectangle (6,4);
\draw[accent,line width=1.1pt] (-5.0,0.0)--(-3.0,3.0)--(0.0,-2.0);
\draw[accent,line width=1.1pt] (0.0,-4.0)--(2.0,1.0)--(5.0,-1.0);
\end{scope}
\fill[accent] (-5.0,0.0) circle (2.2pt);
\fill[accent] (0.0,-2.0) circle (2.2pt);
\filldraw[fill=white,draw=accent,line width=0.9pt] (0.0,-4.0) circle (2.2pt);
\fill[accent] (5.0,-1.0) circle (2.2pt);
\end{tikzpicture}\end{center}

\ex{2} Un depósito tiene 20 L de agua. Se llena durante 5 minutos hasta los 120 L, se mantiene así 10 minutos y después se vacía del todo en 10 minutos. Dibuja la gráfica volumen--tiempo y estúdiala. ¿A qué velocidad se llena y a cuál se vacía?

\ex{2} Con 24 m de valla queremos hacer un corral rectangular. Escribe el área en función de la base $x$, haz una tabla para $x=1,\,3,\,6,\,9,\,11$ y di qué dimensiones dan el área máxima.

\ex{2} Lanzamos una pelota hacia arriba y su altura en metros es $h(t)=30t-5t^2$ ($t$ en segundos). Haz una tabla para $t=0,\,1,\,2,\,3,\,4,\,6$, represéntala y responde: ¿qué altura máxima alcanza? ¿Cuándo vuelve al suelo? ¿Qué dominio tiene sentido?

\ex{3} El gimnasio A cobra 30 € de cuota mensual más 2 € por sesión; el gimnasio B no tiene cuota pero cobra 5 € por sesión. Escribe las dos funciones, represéntalas en los mismos ejes y di a partir de cuántas sesiones al mes conviene el gimnasio A.

\ex{3} La tarifa del agua de un pueblo es 0,50 €/m\textsuperscript{3} hasta 10 m\textsuperscript{3} y 0,80 €/m\textsuperscript{3} por cada m\textsuperscript{3} que pasa de 10. Escribe la función a trozos, calcula el precio de 15 m\textsuperscript{3} y di si es continua.

\ex{3} Ana sale de casa caminando a 5 km/h. Una hora después, su hermano sale por el mismo camino en bicicleta a 15 km/h. Escribe la distancia a casa de cada uno en función del tiempo (desde que sale Ana) y halla cuándo y dónde la alcanza.

\vfill
\begin{center}\color{muted}\small\faCheckCircle[regular]\ \ ¡Ánimo! Si te cuesta algún ejercicio, repasa los apuntes interactivos de la unidad.\end{center}

\end{document}
